% 7.5 磁点群的共表示
\section{磁点群的共表示}

\noindent\textbf{灰群；时间反演不变性的后果}\ 已提到过，时间反演操作的效果是反转原子或离子的磁矩。因此对顺磁或抗磁晶体中的离子，在没有外磁场时，晶体在时间反演操作下不变；于是可以说，普通点群（即 I 型 Shubnikov 点群）不足以描述该离子局部环境的对称性，更好的描述是相应的 II 型 Shubnikov 群之一——灰群。考虑这些群并导出其共表示、进而考察晶体在时间反演操作下的不变性是否引入额外简并，是相当直接了当的（Cracknell 1966\textit{a}, Cracknell and Wong 1967, Dimmock and Wheeler 1962\textit{b}, 1964）。

对 32 个灰群中的每一个，可取 $A$ 为元素 $\theta$（时间反演操作本身）。于是 $\mathbf{G} = \mathbf{M}$、$\mathbf{G}$ 在 $\mathbf{M}$ 中指数为 2，情形属于 (a)、(b) 或 (c) 视 $\Gamma^{i*} \simeq \Gamma^{i}$ 与否：对实表示（(b) 类），时间反演把每个态映到同一能级的简并伙伴（但不产生新的额外简并，除非表示为复的）；对复表示（(c) 类），${}^{1}\Gamma$ 与 ${}^{2}\Gamma$ 这一共轭对在时间反演下互相简并。这就是著名的 Kramers 简并：对含奇数个电子的体系，时间反演不变性保证至少二重简并。

\noindent\textbf{III 型磁点群的共表示}\ 对 58 个黑白点群，逐一选取幺正子群 $\mathbf{G}$、陪集代表 $A$，按情形 (a)/(b)/(c) 由表 2.2 的单值与双值表示导出共表示，并标明时间反演引起的额外简并。原书第 636--650 页逐一列表；以下为扫描。

\begin{figure}[H]
\centering
\includegraphics[width=0.86\textwidth]{c7_p649.jpg}
\caption*{\textbf{7.5 节的共表示表}：磁点群的不可约共表示（原书第 636--650 页扫描）}
\end{figure}
\begin{figure}[H]
\centering
\includegraphics[width=0.86\textwidth]{c7_p650.jpg}
\caption*{原书第 637 页}
\end{figure}
\begin{figure}[H]
\centering
\includegraphics[width=0.86\textwidth]{c7_p651.jpg}
\caption*{原书第 638 页}
\end{figure}
\begin{figure}[H]
\centering
\includegraphics[width=0.86\textwidth]{c7_p652.jpg}
\caption*{原书第 639 页}
\end{figure}
\begin{figure}[H]
\centering
\includegraphics[width=0.86\textwidth]{c7_p653.jpg}
\caption*{原书第 640 页}
\end{figure}
\begin{figure}[H]
\centering
\includegraphics[width=0.86\textwidth]{c7_p654.jpg}
\caption*{原书第 641 页}
\end{figure}
\begin{figure}[H]
\centering
\includegraphics[width=0.86\textwidth]{c7_p655.jpg}
\caption*{原书第 642 页}
\end{figure}
\begin{figure}[H]
\centering
\includegraphics[width=0.86\textwidth]{c7_p656.jpg}
\caption*{原书第 643 页}
\end{figure}
\begin{figure}[H]
\centering
\includegraphics[width=0.86\textwidth]{c7_p657.jpg}
\caption*{原书第 644 页}
\end{figure}
\begin{figure}[H]
\centering
\includegraphics[width=0.86\textwidth]{c7_p658.jpg}
\caption*{原书第 645 页}
\end{figure}
\begin{figure}[H]
\centering
\includegraphics[width=0.86\textwidth]{c7_p659.jpg}
\caption*{原书第 646 页}
\end{figure}
\begin{figure}[H]
\centering
\includegraphics[width=0.86\textwidth]{c7_p660.jpg}
\caption*{原书第 647 页}
\end{figure}
\begin{figure}[H]
\centering
\includegraphics[width=0.86\textwidth]{c7_p661.jpg}
\caption*{原书第 648 页}
\end{figure}
\begin{figure}[H]
\centering
\includegraphics[width=0.86\textwidth]{c7_p662.jpg}
\caption*{原书第 649 页}
\end{figure}
\begin{figure}[H]
\centering
\includegraphics[width=0.86\textwidth]{c7_p663.jpg}
\caption*{原书第 650 页}
\end{figure}

% 7.6 磁空间群的共表示
\section{磁空间群的共表示}

确定磁空间群共表示的问题与确定磁点群共表示的问题之间的关系，非常类似于由点群表示确定 Fedorov 群（普通空间群）表示的问题之间的关系。由点群表示确定普通空间群表示的理论已在 3.7 与 3.8 节考虑过，并在第四章更严格地处理。记得利用第四章的诱导表示理论，空间群表示的确定问题可以简化：对表示域中的每个波矢 $\mathbf{k}$，只需推导小陪群 $\overline{\mathbf{G}}^{\mathbf{k}}$ 的、带有依赖于该 $\mathbf{k}$ 矢量的因子系统的射影表示。类似地，磁空间群的共表示可以通过诱导从小群的共表示得到；II 型（灰）空间群的共表示由相应 Fedorov 群表示的 (a)/(b)/(c) 分类直接给出，III 型（黑白）空间群则需对每个 $\mathbf{k}$ 辨认使其翻转为 $-\mathbf{k}$ 的反幺正元素，并把共表示按三种情形构造。原书第 652--664 页对磁空间群逐一列表；以下为扫描。

\begin{figure}[H]
\centering
\includegraphics[width=0.86\textwidth]{c7_p665.jpg}
\caption*{\textbf{7.6 节的共表示表}：磁空间群的不可约共表示（原书第 652--664 页扫描）}
\end{figure}
\begin{figure}[H]
\centering
\includegraphics[width=0.86\textwidth]{c7_p666.jpg}
\caption*{原书第 653 页}
\end{figure}
\begin{figure}[H]
\centering
\includegraphics[width=0.86\textwidth]{c7_p667.jpg}
\caption*{原书第 654 页}
\end{figure}
\begin{figure}[H]
\centering
\includegraphics[width=0.86\textwidth]{c7_p668.jpg}
\caption*{原书第 655 页}
\end{figure}
\begin{figure}[H]
\centering
\includegraphics[width=0.86\textwidth]{c7_p669.jpg}
\caption*{原书第 656 页}
\end{figure}
\begin{figure}[H]
\centering
\includegraphics[width=0.86\textwidth]{c7_p670.jpg}
\caption*{原书第 657 页}
\end{figure}
\begin{figure}[H]
\centering
\includegraphics[width=0.86\textwidth]{c7_p671.jpg}
\caption*{原书第 658 页}
\end{figure}
\begin{figure}[H]
\centering
\includegraphics[width=0.86\textwidth]{c7_p672.jpg}
\caption*{原书第 659 页}
\end{figure}
\begin{figure}[H]
\centering
\includegraphics[width=0.86\textwidth]{c7_p673.jpg}
\caption*{原书第 660 页}
\end{figure}
\begin{figure}[H]
\centering
\includegraphics[width=0.86\textwidth]{c7_p674.jpg}
\caption*{原书第 661 页}
\end{figure}
\begin{figure}[H]
\centering
\includegraphics[width=0.86\textwidth]{c7_p675.jpg}
\caption*{原书第 662 页}
\end{figure}
\begin{figure}[H]
\centering
\includegraphics[width=0.86\textwidth]{c7_p676.jpg}
\caption*{原书第 663 页}
\end{figure}
\begin{figure}[H]
\centering
\includegraphics[width=0.86\textwidth]{c7_p677.jpg}
\caption*{原书第 664 页}
\end{figure}

% 7.7 例子：P4'2/mnm' 的空间群共表示及其 Kronecker 积
\section{例子：$P4'_{2}/mnm'$ 的空间群共表示及其 Kronecker 积}

作为前几节理论应用的说明，我们考虑推导 III 型 Shubnikov 空间群 $P4'_{2}/mnm'$ 的共表示；这是反铁磁体 $\mathrm{MnF}_{2}$ 的空间群。该空间群曾由 Dimmock and Wheeler (1962\textit{b}) 与 Cracknell (1967\textit{c}) 考虑过。结构示于图 7.5，空间群的元素可由表 7.2 与 3.7 辨认。显然它关联于普通空间群 $P4_{2}/mnm\ (D^{14}_{4h})$；由表 7.2 可见其生成元与 $P4_{2}/mnm$ 相同，只是 $\{C^{+}_{4z} \mid {\frac{1}{2}}{\frac{1}{2}}{\frac{1}{2}}\}$ 必须乘以 $\theta$。处理 $\mathrm{MnF}_{2}$ 的结构时，用《国际表》给出的原点与轴向设定是方便的。原书第 666--689 页逐一推导了各对称点的共表示（按 (a)/(b)/(c) 情形）以及它们的 Kronecker 积约化；以下为扫描。

\begin{figure}[H]
\centering
\includegraphics[width=0.86\textwidth]{c7_p679.jpg}
\caption*{\textbf{7.7 节的表}：$P4'_{2}/mnm'$ 的共表示及其直积（原书第 666--689 页扫描）}
\end{figure}
\begin{figure}[H]
\centering
\includegraphics[width=0.86\textwidth]{c7_p680.jpg}
\caption*{原书第 667 页}
\end{figure}
\begin{figure}[H]
\centering
\includegraphics[width=0.86\textwidth]{c7_p681.jpg}
\caption*{原书第 668 页}
\end{figure}
\begin{figure}[H]
\centering
\includegraphics[width=0.86\textwidth]{c7_p682.jpg}
\caption*{原书第 669 页}
\end{figure}
\begin{figure}[H]
\centering
\includegraphics[width=0.86\textwidth]{c7_p683.jpg}
\caption*{原书第 670 页}
\end{figure}
\begin{figure}[H]
\centering
\includegraphics[width=0.86\textwidth]{c7_p684.jpg}
\caption*{原书第 671 页}
\end{figure}
\begin{figure}[H]
\centering
\includegraphics[width=0.86\textwidth]{c7_p685.jpg}
\caption*{原书第 672 页}
\end{figure}
\begin{figure}[H]
\centering
\includegraphics[width=0.86\textwidth]{c7_p686.jpg}
\caption*{原书第 673 页}
\end{figure}
\begin{figure}[H]
\centering
\includegraphics[width=0.86\textwidth]{c7_p687.jpg}
\caption*{原书第 674 页}
\end{figure}
\begin{figure}[H]
\centering
\includegraphics[width=0.86\textwidth]{c7_p688.jpg}
\caption*{原书第 675 页}
\end{figure}
\begin{figure}[H]
\centering
\includegraphics[width=0.86\textwidth]{c7_p689.jpg}
\caption*{原书第 676 页}
\end{figure}
\begin{figure}[H]
\centering
\includegraphics[width=0.86\textwidth]{c7_p690.jpg}
\caption*{原书第 677 页}
\end{figure}
\begin{figure}[H]
\centering
\includegraphics[width=0.86\textwidth]{c7_p691.jpg}
\caption*{原书第 678 页}
\end{figure}
\begin{figure}[H]
\centering
\includegraphics[width=0.86\textwidth]{c7_p692.jpg}
\caption*{原书第 679 页}
\end{figure}
\begin{figure}[H]
\centering
\includegraphics[width=0.86\textwidth]{c7_p693.jpg}
\caption*{原书第 680 页}
\end{figure}
\begin{figure}[H]
\centering
\includegraphics[width=0.86\textwidth]{c7_p694.jpg}
\caption*{原书第 681 页}
\end{figure}

% 7.8 自旋空间群
\section{自旋空间群}

本章描述的磁点群或磁空间群的概念，基于这样一个假设：算符同时作用于粒子的空间坐标与自旋坐标。这样的操作若要是晶体的对称操作，其结果是原子位置与自旋从初始值起不可分辨。然而，若对磁晶体的自旋哈密顿量的形式作某些简化假定，它可能具有比磁空间群所描述的更多的对称性。于是可以定义并使用\textit{自旋空间群}（spin-space group）$\mathbf{G}_{s}$，它比晶体的磁空间群具有相当多的对称性（Brinkman 1967, Brinkman and Elliott 1966\textit{a}, \textit{b}）。哈密顿量精确形式的细节、连同所涉近似之物理意义的考察，由 Brinkman and Elliott 给出。例如自旋波色散关系的总体特征随后便可确定。自旋空间群是空间群（作用于空间坐标）与自旋转动群（作用于自旋坐标）的半直积子群；其不可约表示可由两个因子的表示构造，且自旋波分支可按其标记。

% 7.9 多色点群与空间群
\section{多色点群与空间群}

Shubnikov 群（黑白群）是通过在位置坐标 $x, y, z$ 之外添加一个取两个可能值的额外坐标 $s$ 而导出的。若让额外坐标 $s$ 不再是二值而是三值、四值或六值的，这一做法可以进一步推广。正如二色点群与空间群的推导，可以推导三色、四色或六色点群与空间群。一般地考虑 $p$ 色点群，$p = 2, 3, 4$ 或 6。设我们有按某种顺序排列的 $p$ 种颜色，例如三色红、绿、黄，则可以定义换色操作 $\mathcal{R}_{p}$，它把这 $p$ 种颜色循环置换，例如 $\mathcal{R}_{3}$ 把红变为绿、绿变为黄、黄变为红（见图 7.10）。$p$ 色点群将有三种类型。I 型 $p$ 色点群（或空间群）就是普通的点群（或空间群）$\mathbf{G}$。II 型 $p$ 色群含有 $\mathcal{R}_{p}$ 的全部幂（至 $p$ 次），从而是灰群对 $p$ 色的推广；III 型 $p$ 色群则含 $\mathcal{R}_{p}$ 的部分幂与 $\mathbf{G}$ 的部分元素之积。$p$ 色群的数目可以由群论论证推出；对 $p = 3, 4, 6$ 的清单与讨论见原书本节（含文献）。多色群在磁结构（多子格磁性）与振动光谱等问题中已有应用。
